% Demo poster for uzhposter.sty: one A0 portrait page that shows every
% building block of the layer -- header, panels, blocks, colours, DAG styles,
% R code, tables and the footer band. Build it here with `make`; a real
% poster lives in its own directory with a two-line Makefile pointing at this
% folder (ds_phd_ws/2026/ is a working example).
\documentclass[10pt]{article}
\usepackage{uzhposter} % resolved from template/latex/poster via TEXINPUTS
\usepackage{amsmath,amssymb}

% A colour swatch for the palette panel: a filled box with the name below it
\newcommand{\swatch}[1]{%
  \begin{minipage}[t]{0.15\linewidth}
    \centering
    {\setlength{\fboxsep}{0pt}\colorbox{#1}{\rule{0pt}{22mm}\hspace{\linewidth}}}\\[2mm]
    \footnotesize\texttt{#1}%
  \end{minipage}}

\begin{document}

\posterheader
  {Department of XYZ}
  {An A0 Poster Layer for the University of Zurich}
  {Martha Muster\textsuperscript{1}, Hans Muster\textsuperscript{1}}
  {\textsuperscript{1}Department of XYZ, University of Zurich ---
   the header takes unit, title, authors and affiliation; the logo comes
   from the shared \texttt{img/} directory}

\begin{posterbody}

% ============================ panel 1 ======================================
\postersection{What this layer gives you}

\begin{minipage}[c]{0.60\linewidth}
  \texttt{uzhposter.sty} turns a plain \texttt{article} document into an A0
  portrait poster in the UZH corporate design. It is print-safe by
  construction: no overlays, no opacity, every element at full strength.

  \vspace{4mm}
  \begin{itemize}
    \setlength{\itemsep}{2mm}
    \item \textbf{Page and type}: A0 portrait geometry, Source Sans Pro, and
      the standard size commands (\texttt{\textbackslash small} \ldots\
      \texttt{\textbackslash Huge}) redeclared for poster reading distance
    \item \textbf{Skeleton}: \texttt{\textbackslash posterheader},
      \texttt{posterbody}, \texttt{\textbackslash postersection},
      \texttt{\textbackslash posterfooter} --- the page you are reading
    \item \textbf{Blocks}: \texttt{posterblock} boxes in the I-Math grey,
      with \texttt{tcolorbox} overrides for emphasis (panel below)
    \item \textbf{Figures}: the \texttt{uzhdag} TikZ styles at poster scale,
      an \texttt{Rcode} verbatim environment, and \texttt{\textbackslash
      fillimage} for cropped photos (right)
    \item \textbf{Helpers} from \texttt{uzhsetup}, e.g.\ the
      DM\uzhsup{3}L superscript geometry via
      \texttt{\textbackslash uzhsup}
  \end{itemize}
\end{minipage}\hfill
\begin{minipage}[c]{0.36\linewidth}
  \begin{center}
    \fillimage[alt={A lecture hall at UZH}]{\linewidth}{200mm}{img/uzh-city.jpg}

    \vspace{2mm}
    {\footnotesize \texttt{\textbackslash fillimage}: cropped to exactly
     this width and height, never stretched}
  \end{center}
\end{minipage}

\vspace{12mm}

% ============================ panel 2 ======================================
\postersection{Building blocks}

% The \vspace{0pt} after each [t] minipage pins its top to the row: a
% minipage[t] otherwise aligns on the baseline of its first box, and a
% tcolorbox carries its baseline at the bottom -- the three columns would
% stagger.
\begin{minipage}[t]{0.31\linewidth}\vspace{0pt}
  \begin{posterblock}{A plain posterblock}
    Grey title bar, almost white body. Math and tables scale with the
    poster fonts:
    \[
      g_j\bigl(\mathrm{E}[X_j \mid \mathbf{Pa}_j]\bigr)
      \;=\; \beta_{0j} + \sum_{X_i \in \mathbf{Pa}_j} \beta_{ij}\, X_i
    \]
    \begin{center}
      \small
      \begin{tabular}{lcc}
        \toprule
        Method & Accuracy & Runtime \\
        \midrule
        Baseline      & 0.81 & 14.2\,ms \\
        \textbf{Ours} & \textbf{0.95} & \textbf{9.8\,ms} \\
        \bottomrule
      \end{tabular}
    \end{center}
  \end{posterblock}
\end{minipage}\hfill
\begin{minipage}[t]{0.31\linewidth}\vspace{0pt}
  \begin{posterblock}{DAG styles at poster scale}
    Same names as on the slides (\texttt{dagnode}, \texttt{dagedge},
    \texttt{dagabsent}), so slide TikZ code pastes over:
    \begin{center}
    \begin{tikzpicture}[x=4.8cm, y=4.0cm]
      \node[dagnode] (am) at (0,0)  {$a_{t-1}$};
      \node[dagnode] (bm) at (0,-1) {$b_{t-1}$};
      \node[dagnode] (a0) at (1,0)  {$a_{t}$};
      \node[dagnode] (b0) at (1,-1) {$b_{t}$};
      \draw[dagedge]  (am) -- (a0);
      \draw[dagedge]  (bm) -- (b0);
      \draw[dagedge, line width=3.6pt] (am) -- (bm);
      \draw[dagabsent] (am) -- (b0);
      \draw[dagabsent] (bm) -- (a0);
    \end{tikzpicture}
    \end{center}
    {\footnotesize dashed berry $=$ absent edge; extra options compose,
     e.g.\ \texttt{line width} for a strong coefficient}
  \end{posterblock}
\end{minipage}\hfill
\begin{minipage}[t]{0.31\linewidth}\vspace{0pt}
  \begin{posterblock}{R code and console output}
    The \texttt{Rcode} environment frames verbatim text in the I-Math grey:
    \begin{Rcode}
fit <- fitAbn(object = mp,
              method = "mle")
> fit$coef$c
  c|intercept          b
     2.417129 -0.3101488
    \end{Rcode}
  \end{posterblock}

  \vspace{6mm}
  \begin{posterblock}[colbacktitle=uzhgreyD, coltitle=white,
      colframe=uzhgreyD]{An emphasis block}
    The optional argument takes \texttt{tcolorbox} overrides --- here a dark
    title bar for the take-home message.
  \end{posterblock}
\end{minipage}

\vspace{12mm}

% ============================ panel 3 ======================================
\postersection{The UZH palette}

The corporate colours carry plain names (article documents cannot easily use
the \texttt{uzh@\ldots} names of the beamer theme); \texttt{uzhimathgrey} is
the grey of the rules, bars and the footer band:

\vspace{5mm}
\begin{center}
  \swatch{uzhblue}\hfill \swatch{uzhcyan}\hfill \swatch{uzhapple}\hfill
  \swatch{uzhgold}\hfill \swatch{uzhorange}\hfill \swatch{uzhberry}
\end{center}

\vspace{12mm}

% ============================ panel 4 ======================================
\postersection{Start your own poster}

\begin{minipage}[t]{0.47\linewidth}\vspace{0pt}
  \begin{posterblock}{1. The two-line Makefile}
    Create a directory with just your \texttt{.tex} and this
    \texttt{Makefile}; nothing from the template is copied:
    \begin{Rcode}
TEMPLATE := ../../template/latex/poster
include $(TEMPLATE)/Makefile
    \end{Rcode}
    \texttt{make} builds the poster PDF, \texttt{make watch} rebuilds on
    every save, \texttt{make clean} drops the artifacts. The beamer-only
    targets (\texttt{draft}, \texttt{short}, \texttt{frames}) are off.
  \end{posterblock}
\end{minipage}\hfill
\begin{minipage}[t]{0.47\linewidth}\vspace{0pt}
  \begin{posterblock}{2. The document skeleton}
    \begin{Rcode}
\documentclass[10pt]{article}
\usepackage{uzhposter}
\begin{document}
\posterheader{Unit}{Title}{Authors}{Affiliation}
\begin{posterbody}
  \postersection{First panel}
  \begin{posterblock}{A block} ... \end{posterblock}
\end{posterbody}
\posterfooter{Reference one \ |\ Reference two}
\end{document}
    \end{Rcode}
  \end{posterblock}
\end{minipage}

\end{posterbody}

\posterfooter{%
  \texttt{\textbackslash posterfooter}: the reference band, bleeding to the
  page edges.
  \ |\
  Muster, M.\ and Muster, H.\ (2026). A poster layer for UZH.
  \emph{J.\ of Reproducible Posters} 1(1).
  \ |\
  Built with \texttt{make} in \texttt{template/latex/poster/}.}

\end{document}
